\section{Architecture}
\label{sec:architecture}

\begin{figure}[t]
\centering
\resizebox{\textwidth}{!}{%
\begin{tikzpicture}[
  font=\small,
  box/.style={draw, rounded corners=2pt, align=center, minimum height=1.05cm, text width=2.5cm, fill=white},
  ext/.style={box, fill=black!5},
  impl/.style={-{Stealth[length=2mm]}, thick},
  spec/.style={-{Stealth[length=2mm]}, thick, dashed},
  lbl/.style={font=\scriptsize, align=center, fill=white, inner sep=1pt}
]
\node[box] (client) at (0,0) {Buyer client\\\scriptsize web app, desktop shell};
\node[box] (bnode) at (0,-2.6) {Buyer node\\\scriptsize keys, vault, wallet};
\node[box] (vnode) at (9.0,-2.6) {Vendor node\\\scriptsize onion mailbox, wallet};
\node[box] (vapp) at (9.0,0) {Vendor client\\\scriptsize opens letters};
\node[ext] (tor) at (4.5,-2.6) {Tor network};
\node[ext] (rep) at (4.5,-0.3) {Replicas and seeds};
\node[ext] (mod) at (13.6,0.4) {Moderator\\\scriptsize dispute only};
\node[ext] (esplora) at (0,-5.0) {Esplora server};
\node[ext] (chain) at (4.5,-5.0) {Bitcoin network};
\node[ext] (rates) at (-4.5,-2.6) {Exchange-rate\\endpoint};

\draw[impl] (client) -- node[lbl, left, xshift=-2pt]{loopback HTTP\\pairing token} (bnode);
\draw[spec] (bnode) -- node[lbl, above=3pt]{sealed\\letter} (tor);
\draw[spec] (tor) -- node[lbl, above=3pt]{onion\\service} (vnode);
\draw[impl] (vnode) -- node[lbl, right]{loopback} (vapp);
\draw[spec] (bnode.north east) -- node[lbl, sloped, above]{catalog} (rep.south west);
\draw[spec] (vnode.north west) -- node[lbl, sloped, above]{publish} (rep.south east);
\draw[impl] (bnode) -- node[lbl, left]{addresses, txs} (esplora);
\draw[impl] (bnode.south east) -- node[lbl, sloped, above]{broadcast} (chain);
\draw[impl] (vnode.south west) -- node[lbl, sloped, above]{payments} (chain);
\draw[spec] (vnode) -- node[lbl, sloped, above]{dispute letter} (mod);
\draw[impl] (bnode) -- node[lbl, above=3pt]{clear-\\net} (rates);
\end{tikzpicture}}
\caption{Components and data paths. Solid arrows exist in code on main or an open pull request; dashed arrows are specified but not implemented. In today's demonstration the buyer node and vendor node are the same process, and the sealed letter is written to that process's own mailbox directory. The phone-to-node pairing is not drawn; the client currently reaches only a node on the same machine.}
\label{fig:architecture}
\end{figure}


Figure~\ref{fig:architecture} shows the components. A user runs a client and a node. The node holds the keys, the wallet, the encrypted vault, the verified catalog, and the escrow state; it watches the chain and talks to other nodes. The client renders the catalog, the cart, and the order screens, and holds no keys.

\subsection{The client-node boundary}

A node listens only on a loopback address and refuses requests whose TCP peer is not a loopback address. On first start it creates a 256-bit random pairing token, stored in a file readable only by its owner. Requests that place orders, read or write the address book, or read the buyer's order list must carry the token. Cross-origin responses name only the client's origins, so a web page from another origin in the user's browser cannot read the token or call the node with it. A phone that pairs with a node on another machine needs an authenticated, encrypted channel to it in addition to the token; we leave that pairing transport to future work (\S\ref{sec:future}).

This boundary protects against remote callers and against other web origins in the user's browser. It does not protect against another process running as the same user on the same machine, which can read the token file; we treat that as a compromised device (A6).

Nodes leave escrow letters for each other in a drop that takes no pairing token, since the sender is another node, not the owner's client. The drop only stores the bytes. A letter has effect only if it opens under one of the node's keys and parses as an escrow message for an escrow the node is in, and every message is then checked against the signed terms (\S\ref{sec:escrow}).

The node refuses any client request that names a destination: a field such as \code{peer}, \code{onion}, \code{seed}, \code{replica}, or \code{url}, or any string value containing \code{://} or \code{.onion}. A client therefore cannot direct the node to dial an address of the client's choosing. The vendor is always the store that signed the listing the node already holds, reached at the mailbox that store published.

\subsection{Deployment}

A vendor runs a node that stays awake (a desktop computer or an Android device), publishes the catalog, runs an onion service that accepts sealed letters, and watches the chain. A buyer uses the same software; a buyer on an iPhone pairs the installed web app with their own desktop or Android node. All peer-to-peer traffic, including catalog fetches, publishing, label sets, letter delivery, and escrow messages, runs over Tor between nodes. A moderator is a node like any other, with a published record and an onion mailbox; it needs to be online only to answer disputes.
